\chapter*{Bibliography}
\addcontentsline{toc}{chapter}{Bibliography}
Bibliography Amirouche, F. M. L. 2006. Fundamentals of Multibody Dynamics: Theory and Applications. Boston: Birkh¨auser. Angeles, J. 2003. Fundamentals of Robotic Mechanical Systems (2nd ed.) New York: Springer-Verlag. Armstrong, W. W. 1979. Recursive Solution to the Equations of Motion of an n-Link Manipulator. Proc. 5th World Congress on Theory of Machines and Mechanisms, (Montreal), pp. 1343–1346, July. Armstrong-H´elouvry, B., Dupont, P., and Canudas de Wit, C. 1994. A Survey of Models, Analysis Tools and Compensation Methods for the Control of Machines with Friction. Automatica, vol. 30, no. 7, pp. 1083–1138. Ascher, U. M., Pai, D. K., and Cloutier, B. P. 1997. Forward Dynamics, Elimination Methods, and Formulation Stiffness in Robot Simulation. Int. J. Robotics Research, vol. 16, no. 6, pp. 749–758. Baker, E. J., and Wohlhart, K. 1996. Motor Calculus: A New Theoretical Device for Mechanics. Graz, Austria: Institute for Mechanics, TU Graz. (Translation of von Mises 1924a and 1924b.) Balafoutis, C. A., Patel, R. V., and Misra, P. 1988. Efficient Modeling and Computation of Manipulator Dynamics Using Orthogonal Cartesian Tensors. IEEE J. Robotics \& Automation, vol. 4, no. 6, pp. 665–676. Balafoutis, C. A., and Patel, R. V. 1989. Efficient Computation of Manipulator Inertia Matrices and the Direct Dynamics Problem. IEEE Trans. Systems, Man, and Cybernetics, vol. 19, no. 5, pp. 1313–1321. Balafoutis, C. A., and Patel, R. V. 1991. Dynamic Analysis of Robot Manipulators: A Cartesian Tensor Approach. Boston: Kluwer Academic Publishers. Ball, R. S. 1900. A Treatise on the Theory of Screws. London: Cambridge Univ. Press. Republished 1998. Baraff, D. 1996. Linear-Time Dynamics using Lagrange Multipliers. Proc. ACM SIGGRAPH ’96, New Orleans, August 4–9, pp. 137–146.

Baumgarte, J. 1972. Stabilization of Constraints and Integrals of Motion in Dynamical Systems. Computer Methods in Applied Mechanics and Engineering, vol. 1, pp. 1–16. Bobrow, J. E. 1989. A Direct Minimization Approach for Obtaining the Distance Between Convex Polyhedra. Int. J. Robotics Research, vol. 8, no. 3, pp. 65–76. Brach, R. M. 1991. Mechanical Impact Dynamics. New York: Wiley. Brady, M., Hollerbach, J. M., Johnson, T. L., Lozano-Perez, T., and Mason, M. T. (eds) 1982. Robot Motion: Planning and Control. Cambridge, MA: The MIT Press. Brand, L. 1953. Vector and Tensor Analysis (4th ed.) New York/London: Wiley/Chapman and Hall. Brandl, H., Johanni, R., and Otter, M. 1988. A Very Efficient Algorithm for the Simulation of Robots and Similar Multibody Systems Without Inversion of the Mass Matrix. In Theory of Robots, P. Kopacek, I. Troch \& K. Desoyer (eds.), Oxford: Pergamon Press, pp. 95–100. Cottle, R. W., and Dantzig, G. B. 1968. Complementary Pivot Theory of Mathematical Programming. Linear Algebra and its Applications, vol. 1, pp. 103– 125. Cottle, R. W., Pang, J.-S., and Stone, R. E. 1992. The Linear Complementarity Problem. Boston: Academic Press. Coutinho, M. G. 2001. Dynamic Simulations of Multibody Systems. New York: Springer. Craig, J. J. 2005. Introduction to Robotics: Mechanics and Control (3ed). Upper Saddle River, NJ: Pearson Prentice Hall. Denavit, J., and Hartenberg, R. S. 1955. A Kinematic Notation for Lower Pair Mechanisms Based on Matrices. Trans. ASME J. Applied Mechanics, vol. 22, pp. 215–221. Dubowsky, S., and Papadopoulos, E. 1993. The Kinematics, Dynamics, and Control of Free-Flying and Free-Floating Space Robotic Systems. IEEE Trans. Robotics \& Automation, vol. 9, no. 5, pp. 531–543. Duff, I.S., Erisman, A. M., and Reid, J. K. 1986. Direct Methods for Sparse Matrices. Oxford: Clarendon Press. Duffy, J. 1990. The Fallacy of Modern Hybrid Control Theory that is Based on “Orthogonal Complements” of Twist and Wrench Spaces. J. Robotic Systems, vol. 7, no. 2, pp. 139–144.

Dupont, P., Hayward, V., Armstrong, B., and Altpeter, F. 2002. Single State Elastoplastic Friction Models. IEEE Trans. Automatic Control, vol. 47, no. 5, pp. 787–792. Fang, A. C., and Pollard, N. S. 2003. Efficient Synthesis of Physically Valid Human Motion. ACM Trans. Graphics (SIGGRAPH 2003), vol. 22, no. 3, pp. 417–426. Featherstone, R. 1983. The Calculation of Robot Dynamics Using ArticulatedBody Inertias. Int. J. Robotics Research, vol. 2, no. 1, pp. 13–30. Featherstone, R. 1984. Robot Dynamics Algorithms. Ph. D. Thesis, Edinburgh University. Featherstone, R. 1987. Robot Dynamics Algorithms. Boston: Kluwer Academic Publishers. Featherstone, R. 1999a. A Divide-and-Conquer Articulated-Body Algorithm for Parallel O(log(n)) Calculation of Rigid-Body Dynamics. Part 1: Basic Algorithm. Int. J. Robotics Research, vol. 18, no. 9, pp. 867–875. Featherstone, R. 1999b. A Divide-and-Conquer Articulated-Body Algorithm for Parallel O(log(n)) Calculation of Rigid-Body Dynamics. Part 2: Trees, Loops and Accuracy. Int. J. Robotics Research, vol. 18, no. 9, pp. 876–892. Featherstone, R. 2004. An Empirical Study of the Joint Space Inertia Matrix. Int. J. Robotics Research, vol. 23, no. 9, pp. 859–871. Featherstone, R. 2005. Efficient Factorization of the Joint Space Inertia Matrix for Branched Kinematic Trees. Int. J. Robotics Research, vol. 24, no. 6, pp. 487–500. Gautier, M., and Khalil, W. 1990. Direct Calculation of Minimum Set of Inertial Parameters of Serial Robots. IEEE Trans. Robotics \& Automation, vol. 6, no. 3, pp. 368–373. Gear, C. W. 1971. Numerical Initial Value Problems in Ordinary Differential Equations. Englewood Cliffs, NJ: Prentice-Hall. George, A., and Liu. J. W. H. 1981. Computer Solution of Large Sparse Positive Definite Systems. Englewood Cliffs, NJ: Prentice-Hall. Gilardi, G., and Sharf, I. 2002. Literature Survey of Contact Dynamics Modelling. Mechanism \& Machine Theory, vol. 37, no. 10, pp. 1213–1239. Gilbert, E. G., Johnson, D. W., and Keerthi, S. S. 1988. A Fast Procedure for Computing the Distance Between Complex Objects in Three-Dimensional Space. IEEE J. Robotics \& Automation, vol. 4, no. 2, pp. 193–203.

Gottschalk, S., Lin, M. C., and Manocha, D. 1996. OBBTree: A Hierarchical Structure for Rapid Interference Detection. Proc. ACM SIGGRAPH ’96, New Orleans, August 4–9, pp. 171–180. Greenwood, D. T. 1988. Principles of Dynamics. Englewood Cliffs, NJ: Prentice-Hall. Haessig, D. A., and Friedland, B. 1991. On the Modeling and Simulation of Friction. Trans. ASME J. Dynamic Systems, Measurement \& Control, vol. 113, no. 3, pp. 354–362. He, X., and Goldenberg, A. A. 1989. An Algorithm for Efficient Computation of Dynamics of Robotic Manipulators. Proc. Fourth Int. Conf. Advanced Robotics, (Columbus, OH), pp. 175–188, June. Hollerbach, J. M. 1980. A Recursive Lagrangian Formulation of Manipulator Dynamics and a Comparative Study of Dynamics Formulation Complexity. IEEE Trans. Systems, Man, and Cybernetics, vol. SMC-10, no. 11, pp. 730– 736. Hooker, W. W., and Margulies, G. 1965. The Dynamical Attitude Equations for an n-Body Satellite. J. Astronautical Sciences, vol. 12, no. 4, pp. 123–128. Hooker, W. W. 1970. A Set of r Dynamical Attitude Equations for an Arbitrary n-Body Satellite having r Rotational Degrees of Freedom. AIAA Journal, vol. 8(2), no. 7, pp. 1205–1207. Hu, W., Marhefka, D. W., and Orin, D. E. 2005. Hybrid Kinematic and Dynamic Simulation of Running Machines. IEEE Trans. Robotics, vol. 21, no. 3, pp. 490–497. Hunt, K. H. 1978. Kinematic Geometry of Mechanisms. Oxford: Oxford University Press. Huston, R. L. 1990. Multibody Dynamics. Boston: Butterworths. Jain, A. 1991. Unified Formulation of Dynamics for Serial Rigid Multibody Systems. J. Guidance, Control, and Dynamics, vol. 14, no. 3, pp. 531–542. Jain, A., and Rodriguez, G. 1993. An Analysis of the Kinematics and Dynamics of Underactuated Manipulators. IEEE Trans. Robotics and Automation, vol. 9, no. 4, pp. 411–422. Jimenez, P., Thomas, F., and Torras, C. 2001. 3D Collision Detection: A Survey. Computers \& Graphics-UK, vol. 25, no. 2, pp. 269–285. Kahn, M. E., and Roth, B. 1971. The Near Minimum-time Control of Openloop Articulated Kinematic Chains. J. Dynamic Systems, Measurement, and Control, vol. 93, pp. 164–172.

Khalil, W., and Kleinfinger, J. F. 1987. Minimum Operations and Minimum Parameters of the Dynamic Models of Tree Structure Robots. IEEE J. Robotics \& Automation, vol. 3, no. 6, pp. 517–526. Khalil, W., and Dombre, E. 2002. Modeling, Identification and Control of Robots. New York, NY: Taylor \& Francis. Khatib, O. 1987. A Unified Approach to Motion and Force Control of Robot Manipulators: The Operational Space Formulation. IEEE J. Robotics \& Automation, vol. 3, no. 1, pp. 43–53. Khatib, O. 1995. Inertial Properties in Robotic Manipulation: An Object-Level Framework. Int. J. Robotics Research, vol. 14, no. 1, pp. 19–36. Lathrop, R. H. 1985. Parallelism in Manipulator Dynamics. Int. J. Robotics Research, vol. 4, no. 2, pp. 80–102. Lilly, K. W., and Orin, D. E. 1991. Alternate Formulations for the Manipulator Inertia Matrix. Int. J. Robotics Research, vol. 10, no. 1, pp. 64–74. Lilly, K. W. 1993. Efficient Dynamic Simulation of Robotic Mechanisms. Boston: Kluwer Academic Publishers. L¨otstedt, P. 1981. Coulomb Friction in Two-Dimensional Rigid Body Systems. Zeitschrift f¨ur Angewandte Mathematik und Mechanik, vol. 61, no. 12, pp. 605–615. L¨otstedt, P. 1982. Mechanical Systems of Rigid Bodies Subject to Unilateral Constraints. SIAM J. Applied Mathematics, vol. 42, no. 2, pp. 281–296. L¨otstedt, P. 1984. Numerical Simulation of Time-Dependent Contact and Friction Problems in Rigid Body Mechanics. SIAM J. Scientific and Statistical Computing, vol. 5, no. 2, pp. 370–393. Longman, R. W., Lindberg, R. E., and Zedd, M. F. 1987. Satellite-Mounted Robot Manipulators—New Kinematics and Reaction Moment Compensation. Int. J. Robotics Research, vol. 6, no. 3, pp. 87–103. Luh, J. Y. S., Walker, M. W., and Paul, R. P. C. 1980a. On-Line Computational Scheme for Mechanical Manipulators. Trans. ASME, J. Dynamic Systems, Measurement \& Control, vol. 102, no. 2, pp. 69–76. Luh, J. Y. S., Walker, M. W., and Paul, R. P. C. 1980b. Resolved-Acceleration Control of Mechanical Manipulators. IEEE Trans. Automatic Control, vol. 25, no. 3, pp. 468–474. McMillan, S. 1994. Computational Dynamics for Robotic Systems on Land and Under Water. Ph. D. Thesis, The Ohio State University, Dept. Electrical Engineering.

McMillan, S., and Orin, D. E. 1995. Efficient Computation of Articulated-Body Inertias Using Successive Axial Screws. IEEE Trans. Robotics \& Automation, vol. 11, no. 4, pp. 606–611. McMillan, S., Orin, D. E., and McGhee, R. B. 1995. Efficient Dynamic Simulation of an Underwater Vehicle with a Robotic Manipulator. IEEE Trans. Systems, Man \& Cybernetics, vol. 25, no. 8, pp. 1194–1206. McMillan, S., and Orin, D. E. 1998. Forward Dynamics of Multilegged Vehicles Using the Composite Rigid Body Method. Proc. IEEE Int. Conf. Robotics and Automation, (Leuven, Belgium), pp. 464–470, May. Marhefka, D. W., and Orin, D. E. 1999. A Compliant Contact Model with Nonlinear Damping for Simulation of Robotic Systems. IEEE Trans. Systems, Man \& Cybernetics—Part A, vol. 29, no. 6, pp. 566–572. Mirtich, B. 1998. V-Clip: Fast and Robust Polyhedral Collision Detection. ACM Trans. Graphics, vol. 17, no. 3, pp. 177–208. von Mises, R. 1924a. Motorrechnung, ein neues Hilfsmittel der Mechanik [Motor Calculus: a new Theoretical Device for Mechanics]. Zeitschrift f¨ur Angewandte Mathematik und Mechanik, vol. 4, no. 2, pp. 155–181. (English translation: Baker and Wohlhart 1996.) von Mises, R. 1924b. Anwendungen der Motorrechnung [Applications of Motor Calculus]. Zeitschrift f¨ur Angewandte Mathematik und Mechanik, vol. 4, no. 3, pp. 193–213. (English translation: Baker and Wohlhart 1996.) Moon, F. C. 1998. Applied Dynamics. New York: Wiley. Murray, J. J., and Neuman, C. P. 1984. ARM: An Algebraic Robot Dynamic Modeling Program. Proc. IEEE Int. Conf. Robotics and Automation, Atlanta, Georgia, pp. 103–114, March 13–15. Murray, R. M., Li, Z., and Sastry, S. S. 1994. A Mathematical Introduction to Robotic Manipulation. Boca Raton, FL: CRC Press. Nakamura, Y., and Yamane, K. 2000. Dynamics Computation of StructureVarying Kinematic Chains and Its Application to Human Figures. IEEE Trans. Robotics and Automation, vol. 16, no. 2, pp. 124–134. Neuman, C. P., and Murray, J. J. 1987. Customized Computational Robot Dynamics. J. Robotic Systems, vol. 4, no. 4, pp. 503–526. Orin, D. E., McGhee, R. B., Vukobratovic, M., and Hartoch, G. 1979. Kinematic and Kinetic Analysis of Open-chain Linkages Utilizing Newton-Euler Methods. Mathematical Biosciences, vol. 43, pp. 107–130, February.

Orlandea, N., Chace, M. A., and Calahan, D. A. 1977. A Sparsity-Oriented Approach to the Dynamic Analysis and Design of Mechanical Systems—Part 1. Trans. ASME J. Engineering for Industry, vol. 99, no. 3, pp. 773–779. Park, F. C., Bobrow, J. E., and Ploen, S. R. 1995. A Lie Group Formulation of Robot Dynamics. Int. J. Robotics Research, vol. 14, no. 6, pp. 609–618. Paul, B. 1975. Analytical Dynamics of Mechanisms—A Computer Oriented Overview. Mechanism and Machine Theory, vol. 10, no. 6, pp. 481–507. Pfeiffer, F., and Glocker, C. (eds) 2000. Multibody Dynamics with Unilateral Contacts. Vienna: Springer-Verlag. Pl¨ucker, J. 1866. Fundamental Views Regarding Mechanics. Philosophical Transactions, vol. 156, pp. 361–380. Roberson, R. E., and Schwertassek, R. 1988. Dynamics of Multibody Systems, Berlin/Heidelberg: Springer-Verlag. Rodriguez, G. 1987. Kalman Filtering, Smoothing, and Recursive Robot Arm Forward and Inverse Dynamics. IEEE J. Robotics \& Automation, vol. RA-3, no. 6, pp. 624–639. Rodriguez, G., Jain, A., and Kreutz-Delgado, K. 1991. A Spatial Operator Algebra for Manipulator Modelling and Control. Int. J. Robotics Research, vol. 10, no. 4, pp. 371–381. Rooney, J. 1977. A Survey of Representations of Spatial Rotations About a Fixed Point. Environment and Planning B, vol. 4, no. 2, pp. 185–210. Rosenthal, D. E. 1990. An Order n Formulation for Robotic Systems. J. Astronautical Sciences, vol. 38, no. 4, pp. 511–529. Saha, S. K. 1997. A Decomposition of the Manipulator Inertia Matrix. IEEE Trans. Robotics \& Automation, vol. 13, no. 2, pp. 301–304. Selig, J. M. 1996. Geometrical Methods in Robotics. New York: Springer. Shabana, A. A. 2001. Computational Dynamics (2nd ed.) New York: Wiley. Siciliano, B., and Khatib, O. (eds) 2008. Springer Handbook of Robotics. Berlin: Springer. Stejskal, V., and Val´aˇsek, M. 1996. Kinematics and Dynamics of Machinery. New York: Marcel Dekker. Stepanenko, Y., and Vukobratovic, M. 1976. Dynamics of Articulated Openchain Active Mechanisms. Math. Biosciences, vol. 28, pp. 137–170. Stewart, D. E. 2000. Rigid-Body Dynamics with Friction and Impact. SIAM Review, vol. 42, no. 1, pp. 3–39.

Uicker, J. J. 1965. On the Dynamic Analysis of Spatial Linkages Using 4 by 4 Matrices. PhD thesis, Northwestern University. Uicker, J. J. 1967. Dynamic Force Analysis of Spatial Linkages. Trans. ASME J. Applied Mechanics, vol. 34, pp. 418–424. Umetani, Y., and Yoshida, K. 1989. Resolved Motion Rate Control of Space Manipulators with Generalized Jacobian Matrix. IEEE Trans. Robotics \& Automation, vol. 5, no. 3, pp. 303–314. Vafa, Z., and Dubowsky, S. 1990a. The Kinematics and Dynamics of Space Manipulators: The Virtual Manipulator Approach. Int. J. Robotics Research, vol. 9, no. 4, pp. 3–21. Vafa, Z., and Dubowsky, S. 1990b. On the Dynamics of Space Manipulators Using the Virtual Manipulator, with applications to Path Planning. J. Astronautical Sciences, vol. 38, no. 4, pp. 441–472. Vereshchagin, A. F. 1974. Computer Simulation of the Dynamics of Complicated Mechanisms of Robot Manipulators. Engineering Cybernetics, no. 6, pp. 65–70. Walker, M. W., and Orin, D. E. 1982. Efficient Dynamic Computer Simulation of Robotic Mechanisms. Trans. ASME, J. Dynamic Systems, Measurement \& Control, vol. 104, pp. 205–211. Wittenburg, J. 1977. Dynamics of Systems of Rigid Bodies. Stuttgart: B. G. Teubner. Wittenburg, J., and Wolz, U. 1985. The Program Mesa Verde for Robot Dynamics Simulations. In Robotics Research: The Third International Symposium, O. D. Faugeras and G. Giralt (eds.), Cambridge, MA: The MIT Press, pp. 197–204. Woo, L. S., and Freudenstein, F. 1970. Application of Line Geometry to Theoretical Kinematics and the Kinematic Analysis of Mechanical Systems. J. Mechanisms, vol. 5, no. 3, pp. 417–460. Yamane, K. 2004. Simulating and Generating Motions of Human Figures. Berlin: Springer.
